% ---------------------------------------------------------------------------
% DRAFT 9 — Limitations, Ethics/human-subject statement, licensing.
% Base text reused from acl_source/acl_latex.tex (STRUCTURE.md: "Keep").
% Fixes: copyright overclaim removed, human-subject content added honestly
% (author confirmed annotators were personal contacts -- batchmates and
% parents -- uncompensated volunteers; institutional-review status is
% UNRESOLVED and marked as such, never inferred as "not required"),
% cross-model judge referenced, LLM-use footnote (TMLR FAQ: required,
% page-1 footnote, verbatim policy).
% ---------------------------------------------------------------------------

\section*{Limitations}
\label{sec:limitations}

Several limitations should be noted. The automated four-stage scoring pipeline penalizes
semantically correct predictions that differ in number format or surface phrasing; an early
single-judge pilot (Gemini 2.5 Flash, 874 items, Appendix~\ref{app:judge}) reclassified 78.4\% of
flagged items as correct. None of strict,
judge-only, or judge-augmented scoring is ground truth: strict and the judge's own verdict have
opposing error profiles, and a second, cross-model
judge (Section~\ref{sec:judge}) only partially addresses this, since sharing no model family does not
rule out both judges sharing a bias human adjudicators would not. Human adjudication
(Section~\ref{sec:judge}) confirms this concern is not hypothetical: on the 174 disputed items,
phi4-mini---this paper's primary judge, chosen for full coverage rather than higher accuracy on
disputed cases---is the less accurate of the two judges specifically where they disagree.
DeepSeek-R1-Distill has 14 items in this disputed sample (of its 99 strict REG/NUM/TMP errors);
phi4-mini reclassified 10 of these 14 as correct, and human review disagreed with 7 of those 10. This
is a small, disagreement-selected subsample---not sized to estimate what fraction of DeepSeek's
93.9\% reclassification rate would survive full human review---but it is a genuine limitation of the
cross-model judge design worth stating plainly rather than only reporting the more favourable
control-agreement rate (89.1\%). The benchmark also does not cover Hindi--English code-switched
regulatory text that appears in some official documents.

The human reference study involves a single non-specialist evaluator on 60 medium-and-hard items; it
anchors the harder half of the benchmark but supports no population-level claims about human
performance, and a domain specialist would likely score higher. The IAA study covers 180 of 406
items (44.3\%) across three rounds with CON agreement stable at $\kappa$ = 0.645; agreement on
open-ended tasks reflects the same surface-form scoring sensitivity documented in Appendix~\ref{app:judge}.
Few-shot evaluation was conducted on only four top-performing models (Appendix~\ref{app:fewshot}).

One checkpoint identity rests on inference rather than a surviving primary record: Gemma's local
Ollama tag was re-confirmed as \texttt{gemma3:4b} via the API's own model-metadata endpoint
(4,299,915,632 parameters, family \texttt{gemma3}) after an initial ambiguity in the evaluation
script, and this matches the parameter count already stated in our own methods text---which is
consistent evidence, not fully independent confirmation, since the methods text and the re-run
identity could in principle have been aligned to each other rather than to an external record. We
consider this resolved with reasonable confidence but flag the residual uncertainty rather than
overstating it as settled fact.

A structural limitation shared across all evaluations is context-injection: models receive the
relevant passage directly rather than retrieving it from the full document corpus, testing reading
comprehension over a provided excerpt rather than the retrieval or full-document understanding
real-world compliance workflows also require. The benchmark is a static snapshot of SEBI and RBI
regulation as of early 2026 and will need periodic refresh as the regulatory landscape evolves.

Primary annotation was conducted by a single domain-expert annotator, with the independent human
inter-annotator study above providing external validation; Section~\nameref{sec:humansubjects}
below states who the annotators and evaluator were and what oversight, if any, applied. All
questions and gold-standard answers are in English; the official Hindi-language versions of the same
circulars are not included, limiting applicability to multilingual regulatory analysis. Coverage is
also limited to two regulators---extending to the Insurance Regulatory and Development Authority of
India (IRDAI), the Pension Fund Regulatory and Development Authority (PFRDA), and commodity-segment
regulation would substantially broaden the benchmark's coverage of the Indian financial regulatory
ecosystem.

Finally, IndiaFinBench evaluates short extractive responses rather than longer-form generation; it
does not test whether models can draft regulatory summaries, identify compliance gaps, or produce
legally coherent reasoning chains---capabilities that are arguably more consequential for
legal-technology applications and a natural next step for the benchmark.

We also note a construction limitation surfaced by our own analysis rather than external review: the
difficulty label assigned at annotation time (Section~\ref{sec:models}) predicts empirical
difficulty for only 3 of the 12 evaluated models (Section~\ref{sec:difficulty}), and roughly half of
all items do not discriminate between the systems evaluated here (Section~\ref{sec:effsize}). Both
findings bear on how confidently any single-digit accuracy gap on this benchmark, ours or a future
system's, should be read.

\section*{Ethics Statement}
\label{sec:humansubjects}

\textbf{Source data.} IndiaFinBench is built from primary source documents published by the
Securities and Exchange Board of India (sebi.gov.in) and the Reserve Bank of India (rbi.org.in) for
public access. These are works of an Indian government body; we make no claim that they are in the
public domain or free of copyright, and we do not represent the CC~BY~4.0 license below as covering
the source documents themselves. What is released under CC~BY~4.0 is the authors' own contribution:
the questions, reference answers, difficulty and task-type annotations, evaluation harness, and
model predictions. Every item's provenance (source document and, where applicable, URL) is included
in the release so that claims can be traced back to the original regulatory text.

\textbf{Human participants.} Two human contributions appear in this paper: the second annotator in
the inter-annotator agreement study (Section~\ref{sec:iaa}) and the non-specialist evaluator
(Section~\ref{sec:human}). Both were volunteers from the author's personal and academic network
(fellow students and family members), recruited directly rather than through any institutional
subject pool, informed of the study's purpose before participating, and uncompensated. No
demographic, identifying, or personal information about either participant is collected, stored, or
released; the public dataset contains only their task responses (answers to benchmark items), not
any information about who they are. \textbf{Institutional review status:} this activity has not
been reviewed by an institutional ethics board, and we make no determination here about whether such
review would be required under any particular institution's policy---that determination is left to
readers applying their own institution's standards, and to the author's institution before any
subsequent submission of this or related work. We report this status plainly rather than asserting
an exemption we have not confirmed.

\textbf{Content and use.} The benchmark is designed to evaluate model performance on regulatory
reasoning tasks and does not contain toxic, harmful, or privacy-violating content. We are not aware
of a use of this benchmark, or of the models evaluated against it, that would constitute a
foreseeable harm beyond the ordinary risks of applying any LLM evaluation to a regulated domain: a
practitioner treating any model's output here as a substitute for qualified legal or compliance
advice would be misusing the benchmark's findings, which concern relative model performance on
constructed QA items, not the models' fitness for real regulatory decision-making.

\subsubsection*{Broader Impact Statement}

This work measures how scoring methodology affects reported LLM performance and releases a new
evaluation resource; it does not train or release a new model, and we see no material risk of
negative societal impact beyond the general risks of LLM benchmarking research. To the extent this
paper has a broader implication, it argues for more transparent, multi-regime benchmark reporting,
which we expect to reduce rather than increase the risk of over-trusting a single leaderboard number
in deployment decisions.
